\scanpage{17}
désignant la réduction mod. idéal maximal).

Si $x$ et $y$ sont disjoints, alors \inserted{(et alors seulement)} \deleted{\unsure{l'applicat|l'applicati|l'appl.}}
\[
\begin{aligned}
 (u,v)&\unsure{\mathrel{\vdash\mkern-8mu\dashv}|\mapsto}u\wedge v\\
 D_x\times\unsure{D|D_x|D_y}&\longrightarrow\det V
\end{aligned}
\]
induit
\[
T_x\times T_y\longrightarrow\omega_V
\]
qui se factorise par un \inserted{iso}morphisme
\[
\underbrace{T_x\wedge T_y}_{\substack{\text{produit contracté}\\\text{de }\Lambda^{\unsure{*|\alpha}}\text{-torseurs}}}\simeq\omega_V
\tag*{(5)}
\]
\deleted{\unsure{appelé|noté}} l'image \inserted{dans $\omega_V$} de $(u\wedge v)$ par cette application sera encore noté $u\wedge v$. On aura
\[
u\wedge v=-v\wedge u
\tag*{(6)}
\]
\margin{NB $-1\in\Lambda^*$, d'où \unsure{pour|pose}, pour un $\Lambda^*$-torseur $T$, et $u\in T$, $-u\overset{\mathrm{def}}{=}(-1).u$ )}

Soient $x,y$ \inserted{$\in P(V)^!$} disjoints, et $z\in P(V)^!$ quelc, soient
\[
u\in T_{\unsure{a|x}},\ v\in T_y,\ w\in T_z
\]
considérons la relation
\[
u\wedge v=v\wedge w=w\wedge u
\tag*{(7)}
\]
équivalente à
\[
\deleted{\unsure{u|v}}w\wedge v=v\wedge u=u\wedge w\qquad,
\tag*{(7$'$)}
\]
on vérifie de suite qu'elle équivaut à
\[
u+v+w=0
\tag*{(8)}
\]
et qu'elle implique que $z$ est disjoint de $x,y$.
% archive sheet number: 16
